This section records the formalism choices, simulator capabilities, and
conventions that apply to every experiment in this repository.

% ─────────────────────────────────────────────────────────────────────────────
\subsection*{Foundational Reference}

The tick-counter experiment reproduced in this repository originates from:

\begin{quote}
  Damian Vicino, Olivier Dalle, Gabriel A.\ Wainer.
  \textit{A Data Type for Discretized Time Representation in DEVS.}
  SIMUTOOLS 2014. \texttt{hal-01055555}.
\end{quote}

% ─────────────────────────────────────────────────────────────────────────────
\subsection*{CDBoost PDEVS}

CDBoost implements Parallel DEVS (PDEVS) as originally formulated.
The original PDEVS formalism does not specify a port system; CDBoost
follows this faithfully.  Coupling is expressed as a list of
\texttt{(sender, receiver)} model pointer pairs; all submodels in a
coupled model share a single message type \texttt{MSG}.  This is a
deliberate design choice, not a restriction.

Multiple simultaneous messages are delivered as a \textbf{message bag}
(\texttt{std::vector<MSG>}).

When a model is both imminent and receives external input, the
\textbf{confluence function} $\delta_{\text{con}}$ resolves the two transitions.
CDBoost's default (and the convention used in these experiments) is
internal-first:
\[
  \delta_{\text{con}}(s, e, x) =
    \delta_{\text{ext}}\!\left(\delta_{\text{int}}(s),\; 0,\; x\right)
\]

\subsubsection*{Coupled Model Structure}

\texttt{coupled<TIME>} takes four constructor arguments:
\begin{itemize}
  \item \texttt{models} --- all submodel pointers
  \item \texttt{eic}    --- external input coupling (models that receive external input)
  \item \texttt{ic}     --- internal coupling as \texttt{(sender, receiver)} pairs
  \item \texttt{eoc}    --- external output coupling (models whose output leaves the coupled model)
\end{itemize}

% ─────────────────────────────────────────────────────────────────────────────
\subsection*{Message Type}

Because \texttt{MSG} is shared across a coupled model, semantically
distinct signals must be encoded in the value.  The convention used in
these experiments:

\begin{center}
\begin{tabular}{@{}ll@{}}
  \toprule
  \texttt{MSG} value & Semantics \\
  \midrule
  $v \neq 0$ & add $v$ to the counter \\
  $0$        & output current count, then reset to zero \\
  \bottomrule
\end{tabular}
\end{center}

Concretely, \texttt{MSG = std::any} wrapping \texttt{int}.  The
\texttt{infinite\_counter} basic model in CDBoost expects this encoding.

% ─────────────────────────────────────────────────────────────────────────────
\subsection*{TIME Type}

CDBoost is parameterised over a \texttt{TIME} template argument.
\texttt{time\_inf<TIME>::value()} provides infinity; a specialization is
required for types without \texttt{std::numeric\_limits<TIME>::infinity()}.

Common choices:

\begin{center}
\begin{tabular}{@{}lll@{}}
  \toprule
  Type & Representation & Risk \\
  \midrule
  \texttt{double} & IEEE 754 double & $1/10$ not exact; errors may break causality chain and are unbounded \\
  \texttt{boost::rational<int>} & exact fraction & no rounding error; heavier compute \\
  \bottomrule
\end{tabular}
\end{center}

CDBoost ships \texttt{cdboost::rational\_time.hpp} which specialises
\texttt{time\_inf<boost::rational<int>{}>} so that \texttt{boost::rational}
satisfies the TIME requirements without further wrapping.
